\documentclass[11pt]{article}

\usepackage[margin=1in]{geometry}
\usepackage{amsmath,amssymb,mathtools}
\usepackage{graphicx}
\usepackage{float}
\usepackage{microtype}
\usepackage[hidelinks]{hyperref}

\newcommand{\cF}{\mathcal F}
\newcommand{\cH}{\mathcal H}
\newcommand{\cM}{\mathcal M}
\newcommand{\cV}{\mathcal V}
\newcommand{\PP}{\mathbb P}
\newcommand{\vp}{\varphi}

\title{Sphere Five-Point Virasoro Block in the Pillow Frame\\
and the Sphere \texorpdfstring{$n$}{n}-Point Open-Necklace
Elliptic \texorpdfstring{$h$}{h}-Recursion}
\author{}
\date{August 29, 2026}

\begin{document}
\maketitle

\begin{abstract}
We write the chiral sphere five-point Virasoro block as a matrix element on
the pillow $T^2/\mathbb Z_2$.  Four external operators sit at the branch
points, while the fifth becomes a local operator inserted between two
projected propagation segments.  We first audit the sphere four-point
elliptic recursion directly against a PBW expansion.  This fixes the typo in
Eq.~(E.103) of the string note and separates the geometric
$c$-normalization from the character-absorbed $c-1$ normalization.  We then
generalize the corrected conformal factor to five points and state the
resulting two-channel elliptic $h$-recursion.  The recursion uses the
correlated large-weight pillow-character asymptotic as its regular seed.
Finally, we extend the construction to a sphere $n$-point block with
$n-4$ local insertions during pillow evolution.  The resulting
$(n-3)$-edge recursion has a single diagonal large-weight parameter and the
same twisted pillow character.
\end{abstract}

\section{Sphere block and pillow geometry}

Place the five chiral primaries at
\begin{equation}
 V_{d_1}(0),\qquad V_{d_2}(z),\qquad V_{d_3}(1),\qquad
 V_{d_4}(\infty),\qquad V_{d_5}(t),
 \label{eq:insertions}
\end{equation}
and choose $(0,z,1,\infty)$ as the branch quartet.  In the comb channel the
block is
\begin{equation}
 \cF^{S^2}_5(z,t)
 =
 \langle d_4|V_{d_3}(1)\PP_{h_2}V_{d_5}(t)
 \PP_{h_1}V_{d_2}(z)|d_1\rangle .
 \label{eq:sphere-block}
\end{equation}
Primary three-point constants are stripped from this definition.

The pillow coordinate and elliptic modulus are
\begin{align}
 w(x)
 &=\frac{1}{\theta_3(\tau)^2}
   \int_0^x\frac{dy}{\sqrt{y(1-y)(z-y)}},
 &
 \tau&=i\frac{K(1-z)}{K(z)},
 \label{eq:pillow-map}\\
 q&=e^{i\pi\tau},
 &
 z&=\frac{\theta_2(\tau)^4}{\theta_3(\tau)^4}.
 \label{eq:elliptic-data}
\end{align}
Here $\theta_j(\tau)=\vartheta_j(0|\tau)$.  The fifth puncture maps to
\begin{equation}
 w_5=w(t),\qquad
 \frac{dw_5}{dt}
 =\frac{1}{\theta_3(\tau)^2\sqrt{t(1-t)(z-t)}}.
 \label{eq:mobile-jacobian}
\end{equation}
The square-root branch in Eqs.~\eqref{eq:pillow-map} and
\eqref{eq:mobile-jacobian} must be continued together with all fractional
powers below.

Quantizing the pillow along the $A$-cycle splits its propagation at $w_5$.
The raw cylinder variables are
\begin{equation}
 \widetilde p_1=e^{iw_5},\qquad
 \widetilde p_2=e^{i(\pi\tau-w_5)}.
\end{equation}
After rotating each boundary state by half an $A$-cycle, define the
OPE-aligned local coordinates
\begin{equation}
 p_1=-e^{iw_5},\qquad
 p_2=-e^{i(\pi\tau-w_5)},\qquad
 \boxed{p_1p_2=q}.
 \label{eq:segment-nomes}
\end{equation}
They obey
\begin{equation}
 \frac{z}{t}=4p_1+O(p^2),\qquad t=4p_2+O(p^2),
 \label{eq:aligned-cusp}
\end{equation}
which fixes the signs in the recursion variables $4p_i$.  The half-period
phases are absorbed into the aligned pillow boundary states.
The corresponding projected pillow matrix element is
\begin{equation}
 \boxed{
 \cM_5(q,w_5)=
 \langle\psi_{43}|\PP_{h_2}
 p_2^{L_0-c/24}V_{d_5}^{(w)}(0)
 p_1^{L_0-c/24}\PP_{h_1}|\psi_{12}\rangle .}
 \label{eq:pillow-matrix-element}
\end{equation}
The states $|\psi_{12}\rangle$ and $\langle\psi_{43}|$ are prepared by the
two lower and two upper corner insertions, respectively.  They include the
corner regularizations appropriate to the pillow metric.

\section{Exact conformal factor}

The four branch-point operators produce precisely the Weyl-anomaly and
corner factor of the four-point pillow relation.  Since the fifth operator
is away from the ramification points, it transforms as an ordinary primary,
\begin{equation}
 V_{d_5}^{(x)}(t)
 =\left(\frac{dw_5}{dt}\right)^{d_5}V_{d_5}^{(w)}(w_5).
 \label{eq:primary-transform}
\end{equation}
Consequently, the exact five-point generalization of the four-point
relation is
\begin{align}
 \cF^{S^2}_5(z,t)
 ={}&
 \theta_3(\tau)^{c/2-4(d_1+d_2+d_3+d_4)}
 z^{c/24-d_1-d_2}(1-z)^{c/24-d_2-d_3}
 \left(\frac{dw_5}{dt}\right)^{d_5}
 \cM_5(q,w_5).
 \label{eq:e103-five-point}
\end{align}
The exponent of $1-z$ pairs the operators at $z$ and $1$.  The
$h_3+h_4$ exponent printed in E.103 of the current string-note PDF is a
labeling error, while the $h_2+h_3$ pairing in E.106 is correct.  The
four-point audit in Sec.~\ref{sec:four-point-audit} fixes this statement
independently of the five-point proposal.
It is useful to define, for arbitrary $\kappa$,
\begin{align}
 \Lambda_5^{(\kappa)}(z,t)
 :={}&
 \theta_3(\tau)^{\kappa/2-4(d_1+d_2+d_3+d_4)-2d_5}
 z^{\kappa/24-d_1-d_2}(1-z)^{\kappa/24-d_2-d_3}
 \nonumber\\[-2mm]
 &\hspace{35mm}\times
 [t(1-t)(z-t)]^{-d_5/2}.
 \label{eq:lambda-five}
\end{align}
Then Eq.~\eqref{eq:e103-five-point} becomes
\begin{equation}
 \boxed{
 \cF^{S^2}_5(z,t)=\Lambda_5^{(c)}(z,t)\cM_5(q,w_5).}
 \label{eq:exact-factorization}
\end{equation}

\section{Cusp normalization and large-weight seed}

The factors $4p_1$ and $4p_2$ are fixed by the local degeneration
coordinates; they are not introduced by arbitrarily splitting $16q$.
We first recall how $16q$ itself arises in the four-point problem.  If the
external weights are $\eta_i$ and the exchanged weight is $h$, the plane
OPE normalization gives
\begin{equation}
 \cF^{S^2}_4(z)\sim z^{h-\eta_1-\eta_2}
 \qquad(z\to0).
 \label{eq:four-point-ope-leading}
\end{equation}
The corrected pillow relation instead has the leading form
\begin{equation}
 \cF^{S^2}_4(z)
 \sim z^{c/24-\eta_1-\eta_2}\cM_4(h;q).
 \label{eq:four-point-pillow-leading}
\end{equation}
Matching Eqs.~\eqref{eq:four-point-ope-leading} and
\eqref{eq:four-point-pillow-leading} fixes
\begin{equation}
 \cM_4(h;q)\sim z^{h-c/24}.
 \label{eq:four-point-matrix-leading-z}
\end{equation}
Since the modular lambda function obeys
\begin{equation}
 z=\frac{\theta_2(\tau)^4}{\theta_3(\tau)^4}
  =16q\bigl(1-8q+44q^2+O(q^3)\bigr),
 \label{eq:lambda-leading-sixteen}
\end{equation}
we obtain, without taking any large-weight limit,
\begin{equation}
 \boxed{\cM_4(h;q)\sim(16q)^{h-c/24}.}
 \label{eq:four-point-sixteen-q}
\end{equation}
The cylinder propagator supplies $q^{h-c/24}$; the factor
$16^{h-c/24}$ is fixed by the normalization of the pillow boundary states
relative to the plane OPE coordinate.

For five points the corresponding argument uses the two independent OPE
coordinates
\begin{equation}
 x:=\frac{z}{t}\longrightarrow0,\qquad t\longrightarrow0.
 \label{eq:double-cusp-variables}
\end{equation}
The primary term of the plane comb block is
\begin{equation}
 \cF^{S^2}_5(z,t)
 \sim z^{h_1-d_1-d_2}t^{h_2-h_1-d_5}.
 \label{eq:five-point-ope-leading}
\end{equation}
In the same limit $\theta_3\to1$, $1-z\to1$, and
Eq.~\eqref{eq:mobile-jacobian} gives
\begin{equation}
 \frac{dw_5}{dt}\sim\frac{\epsilon}{t},
 \label{eq:mobile-jacobian-cusp}
\end{equation}
where the constant branch phase $\epsilon$ is absorbed into the local
operator convention.  The exact factorization
Eq.~\eqref{eq:exact-factorization} therefore reduces to
\begin{equation}
 \cF^{S^2}_5(z,t)
 \sim z^{c/24-d_1-d_2}t^{-d_5}\cM_5.
 \label{eq:five-point-pillow-leading}
\end{equation}
Dividing Eq.~\eqref{eq:five-point-ope-leading} by
Eq.~\eqref{eq:five-point-pillow-leading} gives
\begin{align}
 \cM_5
 &\sim z^{h_1-c/24}t^{h_2-h_1}
  =x^{h_1-c/24}t^{h_2-c/24}.
 \label{eq:five-point-matrix-leading-plane}
\end{align}

The inverse pillow map fixes the normalization of each local coordinate
separately:
\begin{align}
 x&=4p_1\left[1-2p_1+2p_2+O(p^2)\right],
 &
 t&=4p_2\left[1+2p_1-2p_2+O(p^2)\right].
 \label{eq:local-coordinate-leading-four}
\end{align}
Consequently the primary propagation factor is fixed a priori as
\begin{equation}
 \boxed{
 \cM_5
 =(4p_1)^{h_1-c/24}(4p_2)^{h_2-c/24}
  \mathcal R_5(h_1,h_2;p_1,p_2),}
 \qquad
 \mathcal R_5=1+O(p_1,p_2)
 \label{eq:five-point-primary-normalization}
\end{equation}
in the double-cusp expansion.  Only after setting $h_1=h_2=h$ and using
$p_1p_2=q$ does this imply
\begin{equation}
 (4p_1)^{h-c/24}(4p_2)^{h-c/24}
 =(16q)^{h-c/24}.
 \label{eq:four-point-reduction-of-segments}
\end{equation}
Thus the five-point factors are derived from the two local coordinates;
their product reproduces the four-point normalization as a consistency
check, rather than serving as its assumed split.

\pagebreak[3]
We now isolate the genuinely additional large-weight proposal.  Set
\begin{equation}
 h_1=H,\qquad h_2=H+a,
 \label{eq:fixed-difference}
\end{equation}
and take $H\to\infty$ at fixed $a,c,d_i,\tau,w_5$.  We assume only that
the regular factor in Eq.~\eqref{eq:five-point-primary-normalization}
approaches the pillow character,
\begin{equation}
 \mathcal R_5(H,H+a;p_1,p_2)
 \longrightarrow\chi_{\mathrm{pill}}(q),
 \qquad
 \chi_{\mathrm{pill}}(q)
 :=\prod_{n=1}^{\infty}(1-q^{2n})^{-1/2}.
 \label{eq:regular-large-h-seed}
\end{equation}
Equivalently, the complete large-weight asymptotic is
\begin{equation}
 \cM_5
 \sim
 (4p_1)^{h_1-c/24}(4p_2)^{h_2-c/24}
 \chi_{\mathrm{pill}}(q).
 \label{eq:large-h-seed}
\end{equation}
The propagation powers in Eq.~\eqref{eq:large-h-seed} are fixed by the
double-cusp normalization.  Only Eq.~\eqref{eq:regular-large-h-seed}, the
limiting form of the regular factor, is the non-geometric assumption
entering the final recursion.

Define the reduced block by
\begin{equation}
 \cM_5
 =
 (4p_1)^{h_1-c/24}(4p_2)^{h_2-c/24}
 \chi_{\mathrm{pill}}(q)
 \cH_5(H,a;p_1,p_2),
 \qquad
 \cH_5\longrightarrow1
 \quad(H\to\infty).
 \label{eq:reduced-block}
\end{equation}
Combining Eqs.~\eqref{eq:exact-factorization} and
\eqref{eq:reduced-block} gives the complete sphere block in the geometric
$c$-normalization:
\begin{equation}
 \boxed{
 \cF^{S^2}_5
 =\Lambda_5^{(c)}(z,t)
 (4p_1)^{h_1-c/24}(4p_2)^{h_2-c/24}
 \chi_{\mathrm{pill}}(q)
 \cH_5(H,a;p_1,p_2).}
 \label{eq:full-block-c}
\end{equation}

The identity
\begin{equation}
 \theta_3(\tau)^{1/2}[z(1-z)]^{1/24}(16q)^{-1/24}
 =\chi_{\mathrm{pill}}(q)^{-1}
 \label{eq:theta-character-identity}
\end{equation}
gives the exactly equivalent Zamolodchikov normalization
\begin{equation}
 \boxed{
 \cF^{S^2}_5
 =\Lambda_5^{(c-1)}(z,t)
 (4p_1)^{h_1-(c-1)/24}(4p_2)^{h_2-(c-1)/24}
 \cH_5(H,a;p_1,p_2).}
 \label{eq:full-block-c-minus-one}
\end{equation}
Thus $\chi_{\mathrm{pill}}$ is explicit in
Eq.~\eqref{eq:full-block-c}, but is already absorbed in
Eq.~\eqref{eq:full-block-c-minus-one}; it must not be included a second
time in the latter convention.

\section{Degenerate data}

We use
\begin{equation}
 c=1+6Q^2,\qquad Q=b+b^{-1},\qquad
 d=\frac{Q^2-\lambda_d^2}{4}.
 \label{eq:momentum-convention}
\end{equation}
The degenerate weight and shifted weight are
\begin{equation}
 h_{r,s}(c)=\frac{Q^2-(rb+s b^{-1})^2}{4},
 \qquad
 h_{r,s}(c)+rs=h_{r,-s}(c),
 \qquad r,s\in\mathbb Z_{>0}.
 \label{eq:kac-weight}
\end{equation}
In this convention the inverse null-norm slope is
\begin{equation}
 A_{r,s}^{c}
 =\frac12
 \prod_{\substack{
 m=1-r,\ldots,r\\
 n=1-s,\ldots,s\\
 (m,n)\neq(0,0),(r,s)}}
 \frac{1}{mb+n b^{-1}}.
 \label{eq:null-slope}
\end{equation}
For weights $d_a,d_b$, define the fusion polynomial
\begin{align}
 P_{r,s}^{c}\!\begin{bmatrix}d_a\\d_b\end{bmatrix}
 :={}&
 \prod_{\substack{p=1-r,\,3-r,\ldots,r-1\\
                   \ell=1-s,\,3-s,\ldots,s-1}}
 \frac{\lambda_{d_a}+\lambda_{d_b}+pb+\ell b^{-1}}{2}
 \frac{\lambda_{d_a}-\lambda_{d_b}+pb+\ell b^{-1}}{2}.
 \label{eq:fusion-polynomial}
\end{align}
Equations~\eqref{eq:null-slope} and \eqref{eq:fusion-polynomial} fix the
normalization of every residue below.

\section{Four-point elliptic recursion and normalization audit}
\label{sec:four-point-audit}

Before proposing the five-point recursion, we fix its normalization by a
direct four-point calculation.  Put four primaries of weights
$(\eta_1,\eta_2,\eta_3,\eta_4)$ at $(0,z,1,\infty)$ and let $h$ propagate
in the $(12)(34)$ channel.  With three-point constants removed, write
\begin{equation}
 \cF^{S^2}_4(z)
 =z^{h-\eta_1-\eta_2}\sum_{N\geq0}F_Nz^N,
 \qquad F_0=1.
 \label{eq:four-point-plane-block}
\end{equation}
The PBW coefficients are obtained without elliptic recursion from
\begin{equation}
 F_N=\rho_{43,N}^{\mathsf T}G_{c,h}^{(N)-1}\rho_{N,12},
 \label{eq:four-point-pbw}
\end{equation}
where $G_{c,h}^{(N)}$ is the level-$N$ Verma Gram matrix.  For example,
\begin{equation}
 F_1=\frac{(h-\eta_1+\eta_2)(h+\eta_3-\eta_4)}{2h}.
 \label{eq:four-point-level-one}
\end{equation}

The corrected geometric pillow relation is
\begin{equation}
 \boxed{
 \cF^{S^2}_4(z)
 =\theta_3(\tau)^{c/2-4\sum_{i=1}^4\eta_i}
  z^{c/24-\eta_1-\eta_2}
  (1-z)^{c/24-\eta_2-\eta_3}\cM_4(h;q).}
 \label{eq:corrected-e103}
\end{equation}
Thus the $(1-z)$ exponent pairs the adjacent insertions at $z$ and $1$.
The $\eta_3+\eta_4$ pairing in E.103 is a typo; the
$\eta_2+\eta_3$ pairing appearing in E.106 is the correct one.

At large $h$ the pillow matrix element has the regular asymptotic
\begin{equation}
 \cM_4(h;q)
 \sim(16q)^{h-c/24}\chi_{\mathrm{pill}}(q),
 \qquad
 \chi_{\mathrm{pill}}(q)
 =\prod_{n\geq1}(1-q^{2n})^{-1/2}.
 \label{eq:four-point-large-h}
\end{equation}
Define the reduced elliptic block by
\begin{equation}
 \cM_4(h;q)
 =(16q)^{h-c/24}\chi_{\mathrm{pill}}(q)\cH_4(h;q).
 \label{eq:four-point-reduced-geometric}
\end{equation}
Using Eq.~\eqref{eq:theta-character-identity} and
$\delta=(c-1)/24$, Eqs.~\eqref{eq:corrected-e103} and
\eqref{eq:four-point-reduced-geometric} become
\begin{equation}
 \boxed{
 \cF^{S^2}_4(z)
 =(16q)^{h-\delta}
  z^{\delta-\eta_1-\eta_2}
  (1-z)^{\delta-\eta_2-\eta_3}
  \theta_3(\tau)^{(c-1)/2-4\sum_i\eta_i}\cH_4(h;q).}
 \label{eq:four-point-zamolodchikov}
\end{equation}
The character is explicit in
Eq.~\eqref{eq:four-point-reduced-geometric}, where the anomaly exponents
contain $c$.  It is already absorbed in
Eq.~\eqref{eq:four-point-zamolodchikov}, where they contain $c-1$.
Consequently, retaining the explicit product in E.106 after making all
$c\to c-1$ shifts counts the same character twice.

The reduced block obeys the standard elliptic recursion
\begin{equation}
\boxed{
 \cH_4(h;q)
 =1+\sum_{r,s\geq1}
 \frac{(16q)^{rs}A_{r,s}^{c}}
      {h-h_{r,s}(c)}
 P_{r,s}^{c}\!\begin{bmatrix}\eta_1\\\eta_2\end{bmatrix}
 P_{r,s}^{c}\!\begin{bmatrix}\eta_4\\\eta_3\end{bmatrix}
 \cH_4\!\left(h_{r,s}(c)+rs;q\right).}
 \label{eq:four-point-h-recursion}
\end{equation}

To decide between the printed formulas numerically, we computed
Eq.~\eqref{eq:four-point-pbw} through level six, changed variables with
$z=\lambda(q)$, and compared every coefficient with
Eq.~\eqref{eq:four-point-h-recursion}.  The three generic test points were
\begin{align*}
 &(26.215;0.13,0.27,0.41,0.56;0.91),\\
 &(31.7;0.19,0.34,0.48,0.67;1.07),\\
 &(42.3;0.16,0.31,0.52,0.73;1.23),
\end{align*}
where each row denotes $(c;\eta_1,\eta_2,\eta_3,\eta_4;h)$.  The largest
relative errors for the corrected normalization were respectively
$1.86\times10^{-10}$, $1.36\times10^{-10}$, and
$7.28\times10^{-11}$.  In contrast, the literal E.103 exponent gives
\begin{equation}
 [q]\bigl(\cH_4^{\mathrm{E.103}}-\cH_4^{\mathrm{rec}}\bigr)
 =16(\eta_2-\eta_4),
 \label{eq:four-point-e103-obstruction}
\end{equation}
while retaining the extra character in the $c-1$ formula gives the
universal obstruction
\begin{equation}
 [q^2]\bigl(\cH_4^{\mathrm{hybrid}}-\cH_4^{\mathrm{rec}}\bigr)
 =-\frac12.
 \label{eq:four-point-character-obstruction}
\end{equation}

The five-point formulas must therefore inherit both lessons: use the
adjacent branch-point pairing $d_2+d_3$, and choose either the geometric
$c$-normalization with $\chi_{\mathrm{pill}}$ or the $c-1$ normalization
without it.  In particular, the reduced block used below is equivalently
extracted as
\begin{equation}
 \boxed{
 \cH_5
 =\frac{\cF^{S^2}_5}
 {\Lambda_5^{(c-1)}(z,t)
  (4p_1)^{h_1-(c-1)/24}(4p_2)^{h_2-(c-1)/24}}.}
 \label{eq:corrected-five-point-extraction}
\end{equation}
No additional character factor is allowed in the denominator of
Eq.~\eqref{eq:corrected-five-point-extraction}.

\section{Elliptic \texorpdfstring{$h$}{h}-recursion for
\texorpdfstring{$\cH_5$}{H5}}

Treat $\cH_5$ as a meromorphic function of $H=h_1$ while holding
$a=h_2-h_1$ fixed.  A pole of the first internal edge occurs at
$H=h_{r,s}$; replacing the null module by its descendant module sends
\begin{equation}
 (H,a)\longmapsto(h_{r,s}+rs,a-rs).
 \label{eq:first-shift}
\end{equation}
A pole of the second edge occurs at $H+a=h_{r,s}$ and sends
\begin{equation}
 (H,a)\longmapsto(h_{r,s}-a,a+rs).
 \label{eq:second-shift}
\end{equation}
The resulting two-channel recursion is
\begin{equation}
\boxed{
\begin{aligned}
 \cH_5(H,a;p_1,p_2)
={}&1\\
&+\sum_{r,s\geq1}
 \frac{(4p_1)^{rs}A_{r,s}^{c}}
      {H-h_{r,s}(c)}
 P_{r,s}^{c}\!\begin{bmatrix}d_1\\d_2\end{bmatrix}
 P_{r,s}^{c}\!\begin{bmatrix}h_{r,s}(c)+a\\d_5\end{bmatrix}
 \cH_5\!\left(h_{r,s}(c)+rs,a-rs;p_1,p_2\right)
 \\
&+\sum_{r,s\geq1}
 \frac{(4p_2)^{rs}A_{r,s}^{c}}
      {H+a-h_{r,s}(c)}
 P_{r,s}^{c}\!\begin{bmatrix}d_4\\d_3\end{bmatrix}
 P_{r,s}^{c}\!\begin{bmatrix}h_{r,s}(c)-a\\d_5\end{bmatrix}
 \cH_5\!\left(h_{r,s}(c)-a,a+rs;p_1,p_2\right).
\end{aligned}}
\label{eq:h-recursion}
\end{equation}
All unshifted arguments $c,d_i,p_1,p_2$ are implicit on the right-hand
side.  Iterating Eq.~\eqref{eq:h-recursion} automatically generates terms
with simultaneous poles in the two internal weights.

When $d_5=0$ and the middle vertex is restricted to the identity channel,
$h_2=h_1$ and $p_1p_2=q$.  The matrix element then reduces to the
four-point pillow propagator, while
\begin{equation}
 (4p_1)^{h-c/24}(4p_2)^{h-c/24}
 =(16q)^{h-c/24},
\end{equation}
providing the normalization check against the four-point result.

\section{Direct PBW comparison}

For an independent check, write the same sphere block in the plane comb
frame as
\begin{equation}
 \cF^{\mathrm{PBW}}_5(z,t)
 =z^{h_1-d_1-d_2}t^{h_2-h_1-d_5}
 \sum_{n_1,n_2\geq0}C_{n_1,n_2}
 \left(\frac{z}{t}\right)^{n_1}t^{n_2},
 \label{eq:pbw-plane-series}
\end{equation}
where each $C_{n_1,n_2}$ is computed directly by contracting the two PBW
Verma bases with inverse Gram matrices and the three descendant
three-point tensors.  No recursion data enter this calculation.

The inverse covering map in the aligned variables has the form
\begin{equation}
 z=16p_1p_2\,Z(p_1p_2),\qquad
 t=4p_2\,Y(p_1,p_2),\qquad
 \frac{z}{t}=4p_1\frac{Z}{Y},
 \label{eq:inverse-covering-series}
\end{equation}
with $Z(0)=Y(0,0)=1$.  Explicitly,
\begin{align}
 Z(q)
 &=\prod_{n\geq1}
   \frac{(1+q^{2n})^8}{(1+q^{2n-1})^8},
 \label{eq:z-cover-product}\\
 Y(p_1,p_2)
 &=(1+p_1)^2
 \prod_{n\geq1}\frac{(1+q^{2n})^4}{(1+q^{2n-1})^4}
 \nonumber\\
 &\quad\times
 \prod_{n\geq1}
 \left[
 \frac{(1+p_1^{2n+1}p_2^{2n})
       (1+p_1^{2n-1}p_2^{2n})}
      {(1+p_1^{2n}p_2^{2n-1})
       (1+p_1^{2n-2}p_2^{2n-1})}
 \right]^2,
 \qquad q=p_1p_2.
 \label{eq:t-cover-product}
\end{align}
Substituting Eqs.~\eqref{eq:inverse-covering-series}--
\eqref{eq:t-cover-product} into Eq.~\eqref{eq:pbw-plane-series} and dividing
by the denominator of Eq.~\eqref{eq:corrected-five-point-extraction} gives
a direct PBW series for $\cH_5$.  Equivalently one may divide by the
geometric factor in Eq.~\eqref{eq:full-block-c}, including exactly one
$\chi_{\mathrm{pill}}$.  The two procedures agree by
Eq.~\eqref{eq:theta-character-identity}.

At first order the direct result is
\begin{equation}
 \cH_5
 =1+
 \frac{2(d_2-d_1)(d_5-a)}{H}\,p_1
 +\frac{2(d_3-d_4)(d_5+a)}{H+a}\,p_2
 +O(p^2),
 \label{eq:first-pbw-coefficients}
\end{equation}
which agrees with the $(r,s)=(1,1)$ terms of
Eq.~\eqref{eq:h-recursion}.  Generic symbolic comparison gives exact
agreement through total bidegree three.  We then computed the PBW Gram
matrices, three-point tensors, and the complete pillow change of variables
with exact rational arithmetic through
\begin{equation}
 n_1+n_2\leq10,
\end{equation}
and compared the result with an independent $80$-digit numerical evaluation
of the $h$-recursion.  This checks all $66$ coefficients in each case.  We
used five fully asymmetric samples:
\begin{center}
\scriptsize
\begin{tabular}{c|c|c|c}
case & $c$ & $(d_1,d_2,d_3,d_4,d_5)$ & $(h_1,h_2)$\\ \hline
1 & $26.215$  & $(0.17,0.29,0.58,0.71,0.43)$ & $(0.93,1.08)$\\
2 & $31.7$    & $(0.21,0.34,0.63,0.79,0.49)$ & $(1.03,1.19)$\\
3 & $42.3$    & $(0.13,0.31,0.69,0.83,0.47)$ & $(0.87,1.14)$\\
4 & $29.3761$ & $(0.113,0.367,1.237,0.524,0.811)$ & $(1.417,0.682)$\\
5 & $37.219$  & $(0.907,0.223,0.376,0.658,1.041)$ & $(0.749,1.533)$
\end{tabular}
\end{center}
The last two samples deliberately use non-monotone external weights, and
case 4 has $a=h_2-h_1=-0.735$; hence both signs of the middle-channel weight
difference are tested.  The discrepancies follow.
\begin{center}
\begin{tabular}{c|c|c|c}
case & $c$ & $(h_1,h_2)$ & largest relative discrepancy\\ \hline
1 & $26.215$ & $(0.93,1.08)$ & $1.363\times10^{-74}$ at $(10,0)$\\
2 & $31.7$   & $(1.03,1.19)$ & $1.225\times10^{-74}$ at $(10,0)$\\
3 & $42.3$   & $(0.87,1.14)$ & $6.019\times10^{-74}$ at $(10,0)$\\
4 & $29.3761$ & $(1.417,0.682)$ & $3.297\times10^{-73}$ at $(10,0)$\\
5 & $37.219$ & $(0.749,1.533)$ & $4.539\times10^{-73}$ at $(0,10)$
\end{tabular}
\end{center}
As a genericity audit, all seven weights are pairwise distinct in each case,
all $65$ nonconstant PBW coefficients through total order ten are nonzero,
and every denominator actually visited by the recursion was recorded.  The
smallest denominator magnitude over all five cases was
$1.232\times10^{-3}$, far from a pole at $80$-digit precision.  Thus the
comparison comprises $5\times66=330$ coefficients and does not rely on a
symmetric choice, a vanishing coefficient, or an accidental null-state
specialization.
Using the literal $h_3+h_4$ exponent printed in E.103 instead produces
\begin{equation}
 [p_1p_2]\bigl(\cH_5^{\mathrm{literal\ E.103}}
                   -\cH_5^{\mathrm{recursion}}\bigr)
 =16(d_2-d_4),
 \label{eq:e103-pbw-obstruction}
\end{equation}
which is nonzero for generic weights.  The direct PBW calculation therefore
both validates Eq.~\eqref{eq:h-recursion} through order ten and fixes the
corrected $1-z$ conformal factor.  As a second regression test, using the
$c-1$ prefactor together with an additional character produces
\begin{equation}
 [p_1^2p_2^2]\bigl(\cH_5^{\mathrm{hybrid}}
                         -\cH_5^{\mathrm{recursion}}\bigr)
 =-\frac12,
 \label{eq:five-point-character-obstruction}
\end{equation}
exactly as predicted by the four-point audit.  This fixes the character
normalization in the five-point computation as well.

\section{Comparison with the CCY plane \texorpdfstring{$h$}{h}-recursion}

There is a second, already-known $h$-recursion for the same five-point
block: the linear-channel plane recursion of Cho--Collier--Yin (CCY),
specialized from their Eq.~(3.26).  It is important not to identify this
representation term by term with Eq.~\eqref{eq:h-recursion}.  Put
\begin{equation}
 x=\frac{z}{t},\qquad H=h_1,\qquad a=h_2-h_1,\qquad
 e_L=d_1-H,\qquad e_R=d_4-H,
 \label{eq:ccy-plane-variables}
\end{equation}
and let $G(H,a,e_L,e_R;x,t)$ be the normalized plane descendant series,
whose regular term is one.  In the notation of this note the CCY recursion
is
\begin{equation}
\boxed{
\begin{aligned}
G(H,a,e_L,e_R;x,t)={}&1\\
&+\sum_{r,s\geq1}\frac{x^{rs}A_{r,s}^{c}}{H-h_{r,s}}
 P_{r,s}^{c}\!\begin{bmatrix}h_{r,s}+e_L\\d_2\end{bmatrix}
 P_{r,s}^{c}\!\begin{bmatrix}h_{r,s}+a\\d_5\end{bmatrix}\\[-2pt]
&\qquad\times
 G(h_{r,s}+rs,a-rs,e_L-rs,e_R-rs;x,t)\\
&+\sum_{r,s\geq1}\frac{t^{rs}A_{r,s}^{c}}{H+a-h_{r,s}}
 P_{r,s}^{c}\!\begin{bmatrix}h_{r,s}-a\\d_5\end{bmatrix}
 P_{r,s}^{c}\!\begin{bmatrix}h_{r,s}-a+e_R\\d_3\end{bmatrix}\\[-2pt]
&\qquad\times G(h_{r,s}-a,a+rs,e_L,e_R;x,t).
\end{aligned}}
\label{eq:ccy-plane-h-recursion}
\end{equation}
Here and below $h_{r,s}=h_{r,s}(c)$.  The pole denominators, the two
shifts of $(H,a)$, and the fusion polynomials adjacent to the middle
insertion agree exactly with Eq.~\eqref{eq:h-recursion}.  The boundary
fusion polynomials and the expansion variables do not: CCY use $(x,t)$
and send $H,d_1,d_4\to\infty$ while keeping $(a,e_L,e_R)$ fixed.  The
pillow limit instead keeps every $d_i$ fixed, uses $(4p_1,4p_2)$, and has
the twisted-character asymptotic before the character is stripped.  Thus
Eq.~\eqref{eq:h-recursion} is not obtained from
Eq.~\eqref{eq:ccy-plane-h-recursion} by the replacements
$x\mapsto4p_1$ and $t\mapsto4p_2$ alone; the coordinate change and the
complete conformal prefactor must be transformed together.

As a normalization and implementation check, we evaluated
Eq.~\eqref{eq:ccy-plane-h-recursion} at $80$-digit precision and compared
all $66$ coefficients with $n_1+n_2\leq10$ against the exact rational PBW
series at the parameters in Eq.~\eqref{eq:value-test-weights} below.  The
largest relative discrepancy is $1.51\times10^{-69}$, at coefficient
$(n_1,n_2)=(10,0)$.  This independently verifies both the CCY dictionary
and the plane-recursion implementation used in the following value test.

\section{Value comparison of the two \texorpdfstring{$h$}{h}-recursions
with the central-charge recursion}

The coefficient audit can be strengthened to a comparison of the block at
finite moduli.  We choose
\begin{equation}
 \begin{gathered}
 c=31.7,\qquad
 (d_1,d_2,d_3,d_4,d_5)=(0.21,0.34,0.63,0.79,0.49),\\
 (h_1,h_2)=(1.03,1.19),
 \end{gathered}
 \label{eq:value-test-weights}
\end{equation}
fix $z=0.08$, and scan $0.12\leq t\leq0.72$.  At each point, the aligned
segment nomes are obtained from the exact product map
\begin{equation}
 p_1p_2=q(z),\qquad t=4p_2Y(p_1,p_2),
 \label{eq:value-test-map}
\end{equation}
not from its leading cusp approximation.  We evaluate the pillow recursion
through total elliptic degree ten and restore the complete plane block with
Eq.~\eqref{eq:full-block-c-minus-one}.  We also evaluate the CCY plane
$h$-recursion, Eq.~\eqref{eq:ccy-plane-h-recursion}, through total degree
ten in $(z/t,t)$.  Independently, the CCY fixed-weight $c$-recursion gives
the same plane descendant series through total degree $22$.

\begin{figure}[H]
 \centering
 \includegraphics[width=0.9\textwidth]{\detokenize{../../Data Set/h-Recursion/sphere_five_pillow_h_vs_c_fixed_z.png}}
 \caption{The reconstructed pillow $h$-recursion, the CCY plane
 $h$-recursion, and the independent CCY $c$-recursion at fixed $z=0.08$.
 The lower panel shows the two relative differences from the order-$22$
 $c$-recursion together with the changes under $8\to10$ in each
 $h$-recursion and $20\to22$ in the $c$-recursion.}
 \label{fig:pillow-h-vs-c-fixed-z}
\end{figure}

For all $61$ sample points, each $h$--$c$ difference is smaller than the sum
of its displayed $h$-recursion shift and the $c$-recursion shift.  For the
pillow representation the median and maximum relative differences are
$1.22\times10^{-9}$ and $2.77\times10^{-5}$; for the CCY plane
$h$-recursion they are $1.81\times10^{-5}$ and $2.52\times10^{-3}$.
At the central point $t=0.42$, the respective differences are
$9.62\times10^{-11}$ and $6.33\times10^{-6}$.  At the least convergent
endpoint $t=0.72$, the order shifts are $2.77\times10^{-7}$ for the pillow
$h$-recursion, $3.20\times10^{-3}$ for the CCY plane $h$-recursion, and
$2.98\times10^{-5}$ for the $c$-recursion.  Therefore both $h$-recursions
converge to the same block within independently measured truncation
uncertainty, while the pillow variables give substantially faster
convergence at equal total order over this scan.

\section{Sphere \texorpdfstring{$n$}{n}-point open-necklace extension}

The five-point construction extends without changing the boundary-state
geometry.  Let
\begin{equation}
 m=n-4,\qquad r=m+1=n-3,
 \label{eq:npoint-counts}
\end{equation}
where $m$ is the number of ordinary mobile insertions and $r$ is the number
of internal propagation segments.  Keep the four branch-point weights
$(d_1,d_2,d_3,d_4)$ and denote the mobile weights by
$\mu_1,\ldots,\mu_m$.  If $w_j=w(t_j)$ are ordered along the chosen pillow
evolution contour, set $w_0=0$, $w_r=\pi\tau$, and define raw segment
parameters
\begin{equation}
 \widetilde p_i=e^{i(w_i-w_{i-1})},\qquad
 \prod_{i=1}^{r}\widetilde p_i=q.
 \label{eq:npoint-raw-segments}
\end{equation}
The phases are fixed by the collision limits of the chosen comb channel.
For the nested real ordering
\begin{equation}
 0<z<t_1<\cdots<t_m<1,
 \qquad
 x_1=\frac{z}{t_1},\quad
 x_i=\frac{t_{i-1}}{t_i}\ (2\leq i\leq m),\quad
 x_r=t_m,
 \label{eq:npoint-plane-ope-variables}
\end{equation}
the simultaneous comb cusp is $x_i\to0$.  The degenerating pillow map gives
\begin{equation}
 e^{iw(t_j)}=-\frac{z}{4t_j}\bigl[1+O(\boldsymbol x)\bigr],
 \qquad q=\frac{z}{16}\bigl[1+O(z)\bigr].
 \label{eq:npoint-pillow-collision-asymptotic}
\end{equation}
It follows directly that, for $r\geq2$,
\begin{align}
 \widetilde p_1
 &=-\frac{x_1}{4}\bigl[1+O(\boldsymbol x)\bigr],
 &
 \widetilde p_i
 &=x_i\bigl[1+O(\boldsymbol x)\bigr]
 &&(2\leq i\leq r-1),
 \nonumber\\
 \widetilde p_r
 &=-\frac{x_r}{4}\bigl[1+O(\boldsymbol x)\bigr].
 \label{eq:npoint-raw-collision-asymptotic}
\end{align}
Therefore the OPE-aligned phases are not arbitrary:
\begin{equation}
 \boxed{
 \varepsilon_1=\varepsilon_r=-1,
 \qquad \varepsilon_i=+1\quad(2\leq i\leq r-1),
 \qquad p_i=\varepsilon_i\widetilde p_i.}
 \label{eq:npoint-aligned-phases}
\end{equation}
Their product is one, and hence
\begin{equation}
 \boxed{\prod_{i=1}^{r}p_i=q.}
 \label{eq:npoint-product-nome}
\end{equation}
For another analytic continuation the phases are continued from
Eq.~\eqref{eq:npoint-aligned-phases}; they are not independently selectable.

Up to translations that place every middle vertex at its sewing seam, the
projected matrix element is
\begin{equation}
\boxed{
 \cM_n=
 \langle\psi_{43}|\PP_{h_r}p_r^{D}V_{\mu_m}
 \PP_{h_{r-1}}p_{r-1}^{D}\cdots
 V_{\mu_1}p_1^{D}\PP_{h_1}|\psi_{12}\rangle,
 \qquad D=L_0-\frac{c}{24}.}
 \label{eq:npoint-pillow-matrix-element}
\end{equation}
All primary three-point constants are stripped.  Since every mobile
operator is inserted away from a branch point, it transforms by its ordinary
primary Jacobian.  In the present conventions,
\begin{equation}
 w(t)=\frac{1}{\theta_3(\tau)^2}
 \int_0^t\frac{dy}{\sqrt{y(1-y)(z-y)}},
 \qquad
 \frac{dw_j}{dt_j}=
 \frac{1}{\theta_3(\tau)^2
 \sqrt{t_j(1-t_j)(z-t_j)}}.
 \label{eq:npoint-pillow-map-jacobian}
\end{equation}

\subsection{Exact sphere-to-pillow conformal factor and extraction recipe}

The relation needed for an actual computation is exact, before taking any
large-internal-weight limit.  It is the $n$-point generalization of the
four-point pillow transformation:
\begin{equation}
 \boxed{\begin{aligned}
 \cF_n^{S^2}(z;\boldsymbol t)
 ={}&
 \theta_3(\tau)^{c/2-4\sum_{a=1}^{4}d_a}
 z^{c/24-d_1-d_2}(1-z)^{c/24-d_2-d_3}
 \\[-1mm]
 &\times
 \prod_{j=1}^{m}\left(\frac{dw_j}{dt_j}\right)^{\mu_j}
 \cM_n(\tau;\boldsymbol w).
 \end{aligned}}
 \label{eq:npoint-exact-sphere-pillow}
\end{equation}
Thus there are no additional pairwise prime-form factors between the
mobile insertions.  Such dependence remains inside the matrix element
$\cM_n$; the coordinate change supplies one local Jacobian per mobile
primary.  All square-root branches are continued together with the pillow
map.

For later use, define for arbitrary $\kappa$
\begin{align}
 \Lambda_n^{(\kappa)}
 :={}&
 \theta_3(\tau)^{\kappa/2-4\sum_{a=1}^{4}d_a
                         -2\sum_{j=1}^{m}\mu_j}
 z^{\kappa/24-d_1-d_2}(1-z)^{\kappa/24-d_2-d_3}
 \nonumber\\[-1mm]
 &\times\prod_{j=1}^{m}
 [t_j(1-t_j)(z-t_j)]^{-\mu_j/2}.
 \label{eq:npoint-lambda}
\end{align}
Equations~\eqref{eq:npoint-pillow-map-jacobian} and
\eqref{eq:npoint-exact-sphere-pillow} are then simply
\begin{equation}
 \boxed{\cF_n^{S^2}=\Lambda_n^{(c)}\cM_n.}
 \label{eq:npoint-exact-lambda-relation}
\end{equation}

Once the effective plumbing coordinates $\rho_i$ are fixed below, the two
equivalent useful forms are
\begin{align}
 \boxed{\cF_n^{S^2}
 ={}
 \Lambda_n^{(c)}
 \left[\prod_i\rho_i^{h_i-c/24}\right]
 \chi_{\mathrm{pill}}(q)\cH_n}
 \label{eq:npoint-full-c}\\
 \boxed{\cF_n^{S^2}
 ={}
 \Lambda_n^{(c-1)}
 \left[\prod_i\rho_i^{h_i-(c-1)/24}\right]\cH_n}
 \label{eq:npoint-full-c-minus-one}
\end{align}
The second line is the most convenient PBW-to-elliptic conversion.  In
particular, the reduced elliptic block is extracted from a direct sphere
block by
\begin{equation}
 \boxed{\begin{aligned}
 \cH_n(H,\boldsymbol a;\boldsymbol p)
 ={}&
 \frac{\cF_n^{S^2}(z;\boldsymbol t)}
 {\Lambda_n^{(c-1)}(z;\boldsymbol t)
  \displaystyle\prod_{i=1}^{r}
  \rho_i^{h_i-(c-1)/24}}
 \\[1mm]
 ={}&
 \frac{\cF_n^{S^2}(z;\boldsymbol t)}
 {\Lambda_n^{(c)}(z;\boldsymbol t)
  \displaystyle\prod_{i=1}^{r}\rho_i^{h_i-c/24}
  \chi_{\mathrm{pill}}(q)}.
 \end{aligned}}
 \label{eq:npoint-extraction-c-minus-one}
\end{equation}
The equality follows from
\begin{equation}
 \theta_3(\tau)^{1/2}[z(1-z)]^{1/24}(16q)^{-1/24}
 =\chi_{\mathrm{pill}}(q)^{-1},
 \label{eq:npoint-theta-character-identity}
\end{equation}
together with $\prod_i\rho_i=16q$.  Notice that all mobile-position and
mobile-weight factors cancel in passing between the two conventions.

The computational recipe is therefore:
\begin{enumerate}
 \item choose the branch quartet $(0,z,1,\infty)$ and order the mobile
 insertions $t_j$ in the desired open-necklace channel;
 \item evaluate $w_j=w(t_j)$ and form the raw segment nomes
 $\widetilde p_i$;
 \item determine the phases from the nested OPE collision limit, then form
 $p_i$ and the endpoint-normalized $\rho_i$;
 \item compute $\Lambda_n^{(c-1)}$ and divide the direct PBW sphere block
 by either denominator in Eq.~\eqref{eq:npoint-extraction-c-minus-one};
 \item re-expand the result in the aligned $p_i$ and apply the edge
 $h$-recursion to the resulting $\cH_n$.
\end{enumerate}

\subsection{Endpoint normalization and the diagonal large-weight seed}

The factors of $4$ belong to the two regulated pillow boundary states, not
to every propagation segment.  The effective plumbing parameters are
therefore
\begin{equation}
 \boxed{\rho_i=4^{\delta_{i1}+\delta_{ir}}p_i,
 \qquad \prod_{i=1}^{r}\rho_i=16q.}
 \label{eq:npoint-effective-plumbing}
\end{equation}
Equations~\eqref{eq:npoint-raw-collision-asymptotic} and
\eqref{eq:npoint-aligned-phases} now give the stronger componentwise
normalization for $r\geq2$
\begin{equation}
 \boxed{\rho_i=x_i\bigl[1+O(\boldsymbol x)\bigr],
 \qquad i=1,\ldots,r.}
 \label{eq:npoint-rho-ope-alignment}
\end{equation}
For $r=1$, define $x_1=z$; then $\rho_1=16q=x_1[1+O(x_1)]$ gives the
same statement.
This formula also covers $r=1$: both Kronecker deltas contribute and give
$\rho_1=16q$, the four-point variable.  For $r=2$ it gives
$(\rho_1,\rho_2)=(4p_1,4p_2)$, while for $r>2$ the intermediate segments
carry no additional factor of $4$.  Assigning $4p_i$ to every segment
would instead give $4^rq$ and would fail the identity-insertion reduction
to the four-point block.

Now set
\begin{equation}
 h_i=H+a_i,\qquad a_1=0,\qquad H\longrightarrow\infty,
 \label{eq:npoint-diagonal-limit}
\end{equation}
with $c,d_a,\mu_j,a_i,p_i$ fixed.  In an orthonormal descendant basis the
$j$th middle vertex obeys the same CCY diagonal limit used at five points,
\begin{equation}
 \frac{\langle H+a_{j+1},A|V_{\mu_j}|H+a_j,B\rangle}
 {\|H+a_{j+1},A\|\,\|H+a_j,B\|}
 \longrightarrow\delta_{AB}.
 \label{eq:npoint-diagonal-vertex}
\end{equation}
Thus all oscillator labels along the open necklace agree.  Their descendant
propagation is
$\prod_i p_i^{|A|}=q^{|A|}$, and the two terminal wavefunctions contract
them exactly as in the four-point pillow amplitude.  Consequently no new
character is generated by an additional middle insertion.  Coefficientwise
in the segment parameters,
\begin{equation}
\boxed{
 \cM_n
 \underset{H\to\infty}{\sim}
 \left[\prod_{i=1}^{r}\rho_i^{h_i-c/24}\right]
 \chi_{\mathrm{pill}}(q),
 \qquad
 \chi_{\mathrm{pill}}(q)=
 \prod_{k\geq1}(1-q^{2k})^{-1/2}.}
 \label{eq:npoint-large-h}
\end{equation}
Equivalently,
\begin{equation}
 \cM_n\sim
 (16q)^{H-c/24}
 \left[\prod_{i=2}^{r}\rho_i^{a_i}\right]
 \chi_{\mathrm{pill}}(q).
 \label{eq:npoint-large-h-fixed-differences}
\end{equation}
The product involving $a_i$ is finite primary propagation; the oscillator
regular part depends only on $q$.  For example, at six points
$r=3$ and $(\rho_1,\rho_2,\rho_3)=(4p_1,p_2,4p_3)$, so this finite factor is
$p_2^{a_2}(4p_3)^{a_3}$ rather than $(4p_2)^{a_2}(4p_3)^{a_3}$.

Define the reduced block by
\begin{equation}
 \cM_n=
 \left[\prod_{i=1}^{r}\rho_i^{h_i-c/24}\right]
 \chi_{\mathrm{pill}}(q)\,
 \cH_n(H,\boldsymbol a;\boldsymbol p),
 \qquad \cH_n\longrightarrow1.
 \label{eq:npoint-reduced-block}
\end{equation}
Substituting this definition into
Eq.~\eqref{eq:npoint-exact-lambda-relation} reproduces the exact full-block
relations~\eqref{eq:npoint-full-c}--\eqref{eq:npoint-full-c-minus-one}.

\subsection{The \texorpdfstring{$(n-3)$}{(n-3)}-edge recursion}

At a pole on edge $i$,
\begin{equation}
 H+a_i=h_{\alpha,\beta}(c).
 \label{eq:npoint-edge-pole}
\end{equation}
The two fusion factors adjacent to that edge are conveniently written
\begin{align}
 \mathcal L_{i;\alpha,\beta}
 &:={}
 \begin{cases}
 P_{\alpha,\beta}^{c}\!\begin{bmatrix}d_1\\d_2\end{bmatrix},&i=1,\\[3pt]
 P_{\alpha,\beta}^{c}\!\begin{bmatrix}
 h_{\alpha,\beta}+a_{i-1}-a_i\\\mu_{i-1}
 \end{bmatrix},&i>1,
 \end{cases}
 \label{eq:npoint-left-residue}\\
 \mathcal R_{i;\alpha,\beta}
 &:={}
 \begin{cases}
 P_{\alpha,\beta}^{c}\!\begin{bmatrix}
 h_{\alpha,\beta}+a_{i+1}-a_i\\\mu_i
 \end{bmatrix},&i<r,\\[3pt]
 P_{\alpha,\beta}^{c}\!\begin{bmatrix}d_4\\d_3\end{bmatrix},&i=r.
 \end{cases}
 \label{eq:npoint-right-residue}
\end{align}
The null-submodule shift is the map
\begin{align}
 \mathsf T_{1;\alpha,\beta}:&\quad
 (H,a_2,\ldots,a_r)\longmapsto
 (h_{\alpha,\beta}+\alpha\beta,
  a_2-\alpha\beta,\ldots,a_r-\alpha\beta),
 \label{eq:npoint-base-shift}\\
 \mathsf T_{i;\alpha,\beta}:&\quad
 (H,a_i)\longmapsto
 (h_{\alpha,\beta}-a_i,a_i+\alpha\beta),
 \qquad i>1,
 \label{eq:npoint-nonbase-shift}
\end{align}
with all $a_{j\ne i}$ fixed in the second line.  These shifts replace the
singular module by the module of weight
$h_{\alpha,\beta}+\alpha\beta$ while preserving every
other physical internal weight at the pole.

The proposed sphere $n$-point pillow recursion is therefore
\begin{equation}
\boxed{
 \cH_n(H,\boldsymbol a;\boldsymbol p)
 =1+
 \sum_{i=1}^{r}\sum_{\alpha,\beta\geq1}
 \frac{\rho_i^{\alpha\beta}A_{\alpha,\beta}^{c}
       \mathcal L_{i;\alpha,\beta}\mathcal R_{i;\alpha,\beta}}
      {H+a_i-h_{\alpha,\beta}(c)}
 \cH_n\!\left(
 \mathsf T_{i;\alpha,\beta}(H,\boldsymbol a);\boldsymbol p\right).}
 \label{eq:npoint-h-recursion}
\end{equation}
For two segments this formula reproduces Eq.~\eqref{eq:h-recursion} term by
term.  For one segment
it is the four-point elliptic recursion.  If all mobile operators are set to
the identity and all $h_i$ are identified, the projectors must first be
composed; Eq.~\eqref{eq:npoint-effective-plumbing} then reduces the entire
propagation to $(16q)^{L_0-c/24}$.  This is a matrix-element reduction and
should not be performed term by term after separating coincident pole
families.

The new part of this proposal is therefore not a new regular function for
each insertion.  Once the fixed-external-weight pillow boundary asymptotic
is accepted, the CCY diagonal-vertex limit forces every additional mobile
operator to act as the identity on the leading oscillator labels.  The same
single twisted character supplies the regular term for all $n$.

\subsection{Six-point PBW consistency certificate}

The first genuine three-edge test is the comb block at
$(0,z,t_1,t_2,1,\infty)$.  Put
\begin{equation}
 x_1=\frac{z}{t_1},\qquad x_2=\frac{t_1}{t_2},\qquad x_3=t_2,
 \qquad q=p_1p_2p_3,
 \label{eq:sixpoint-ope-coordinates}
\end{equation}
and define the two exact mobile-coordinate units by
\begin{equation}
 t_1=4p_2p_3Y_1,\qquad t_2=4p_3Y_2,
 \qquad z=16qZ.
 \label{eq:sixpoint-mobile-units}
\end{equation}
Here $Y_1$ is the five-point pillow product with aggregate nomes
$(p_1,p_2p_3)$, $Y_2$ is the same product with aggregate nomes
$(p_1p_2,p_3)$, and $Z=z/(16q)$.  The direct PBW variables therefore obey
the exact identities
\begin{equation}
 \boxed{
 x_1=4p_1\frac{Z}{Y_1},\qquad
 x_2=p_2\frac{Y_1}{Y_2},\qquad
 x_3=4p_3Y_2.}
 \label{eq:sixpoint-exact-variable-map}
\end{equation}
This equation tests the central assertion that the middle plumbing segment
has no factor of $4$.

The plane primary factor used in the independent PBW computation is
\begin{equation}
 z^{h_1-d_1-d_2}
 t_1^{h_2-h_1-d_3}
 t_2^{h_3-h_2-d_4}.
 \label{eq:sixpoint-plane-primary-factor}
\end{equation}
After multiplying by the inverse of the exact conformal factor in
Eq.~\eqref{eq:npoint-exact-sphere-pillow}, the remaining unit-series factor
is
\begin{align}
 Z^{h_1-c/24}Y_1^{a_2}Y_2^{a_3-a_2}
 &[(1-t_1)(1-z/t_1)]^{d_3/2}
 [(1-t_2)(1-z/t_2)]^{d_4/2}
 \nonumber\\[-1mm]
 &\times
 \theta_3^{-c/2+4(d_1+d_2+d_5+d_6)+2(d_3+d_4)}
 (1-z)^{-c/24+d_2+d_5}\chi_{\mathrm{pill}}^{-1}.
 \label{eq:sixpoint-pbw-stripping-factor}
\end{align}
Together with a factor $4^{n_1+n_3}$ on a PBW monomial
$x_1^{n_1}x_2^{n_2}x_3^{n_3}$, this is the complete direct conversion to
$\cH_6$.

Two independent executable checks were performed.
First, all $20$ coefficients with $n_1+n_2+n_3\leq3$ were computed with
generic symbolic
$(H,a_2,a_3,d_1,\ldots,d_6)$ and
$c=1+6(b+b^{-1})^2$.  Sparse polynomial cross-multiplication proves that
every PBW coefficient equals the corresponding three-edge recursion
coefficient exactly, including all ten coefficients of total degree three.

Second, all $286$ coefficients with total degree at most ten were tested at
each of ten asymmetric generic internal-weight points.  This ensemble
includes strictly ascending and descending triples, central peaks and
valleys, two near-equal nonmonotone triples, and widely separated weights.
The plane PBW Gram matrices,
three-point tensors, coordinate transformation, and conformal-factor
stripping were evaluated with exact rational arithmetic; the recursion was
evaluated independently at $80$ decimal digits.  Before each PBW evaluation,
an exhaustive traversal of every affine recursion shift through total degree
ten rejected any point within $10^{-4}$ of a visited Kac denominator.  The
choices and audit results were
\begin{center}
\scriptsize
\begin{tabular}{c|c|c|c}
case & $(h_1,h_2,h_3)$ & minimum pole denominator & maximum relative error\\
\hline
1 & $(0.9371,1.0837,1.3321)$ & $1.60\times10^{-3}$ & $2.85\times10^{-71}$\\
2 & $(1.4193,0.6871,1.1098)$ & $1.34\times10^{-3}$ & $5.16\times10^{-71}$\\
3 & $(0.7517,1.5372,0.8894)$ & $9.97\times10^{-3}$ & $4.24\times10^{-70}$\\
4 & $(1.6423,1.2711,0.8237)$ & $5.73\times10^{-3}$ & $3.22\times10^{-71}$\\
5 & $(0.9043,0.9187,0.8871)$ & $2.80\times10^{-3}$ & $7.39\times10^{-71}$\\
6 & $(0.6189,1.2473,1.8891)$ & $4.98\times10^{-3}$ & $2.53\times10^{-70}$\\
7 & $(1.1267,1.8429,0.5679)$ & $1.71\times10^{-3}$ & $3.44\times10^{-69}$\\
8 & $(1.7731,0.5427,1.2863)$ & $4.82\times10^{-3}$ & $7.38\times10^{-73}$\\
9 & $(1.1047,1.1299,1.0873)$ & $6.63\times10^{-3}$ & $5.43\times10^{-72}$\\
10 & $(0.4831,1.9637,1.2249)$ & $6.00\times10^{-4}$ & $2.29\times10^{-69}$
\end{tabular}
\end{center}
Cases $1,4,5,10$ use $c=26.215$, cases $2,6,7$ use $c=29.3761$, and
cases $3,8,9$ use $c=37.219$.  Each group reuses one generic external-weight
tuple; the exact tuples are stored in the executable certificate.
All nine weights are distinct at every point, no nonconstant PBW
coefficient vanishes, and no visited recursion denominator is smaller than
the values shown.  Thus $2860$ numerical coefficients pass.  The global
worst case is coefficient $(0,10,0)$ in case 7, precisely a pure excitation
of the unrescaled middle segment.

The executable certificates are the symbolic-order-three and
numerical-order-ten six-point checkers in
\path{Code/sphere_five_kummer_h_recursion/}.

\subsection{Six-point value comparison with the central-charge recursion}

We next compare values of the reconstructed pillow block with the
independent fixed-internal-weight $c$-recursion.  Keep the last two mobile
insertions fixed and vary the first one in the ordered real cell:
\begin{equation}
 t_1=0.32,\qquad t_2=0.62,\qquad 0.015\leq z\leq0.20.
 \label{eq:sixpoint-value-moduli}
\end{equation}
At every $z$, set $q=q(z)$ and let $T(p,q)$ denote the exact five-point
product $4pY(q/p,p)$.  The six-point inverse map is particularly simple:
\begin{equation}
 \boxed{
 T(P_{23},q)=t_1,\qquad T(p_3,q)=t_2,\qquad
 p_1=\frac{q}{P_{23}},\qquad p_2=\frac{P_{23}}{p_3}.}
 \label{eq:sixpoint-value-inverse-map}
\end{equation}
Thus $P_{23}=p_2p_3$ and $p_1p_2p_3=q$ hold exactly; no multivariate root
search is required.  We solved the two monotone real equations using the
infinite product itself.  The largest coordinate residual in the scan was
$7.8\times10^{-67}$.

The generic parameters were
\begin{equation}
 \begin{gathered}
 c=26.215,\qquad
 (d_1,\ldots,d_6)=(0.17,0.29,0.43,0.58,0.71,0.86),\\
 (h_1,h_2,h_3)=(0.9371,1.0837,1.3321),
 \end{gathered}
 \label{eq:sixpoint-value-weights}
\end{equation}
where the external weights are ordered at
$(0,z,t_1,t_2,1,\infty)$.  With
$\delta=(c-1)/24$, the pillow computation was restored to the plane block
as
\begin{equation}
 \cF^{(h,N)}_6=
 \Lambda_6^{(c-1)}
 (4p_1)^{h_1-\delta}p_2^{h_2-\delta}(4p_3)^{h_3-\delta}
 \cH_6^{(N)}(p_1,p_2,p_3).
 \label{eq:sixpoint-value-pillow-reconstruction}
\end{equation}
This is compared with
\begin{equation}
 \cF^{(c,N)}_6=
 z^{h_1-d_1-d_2}t_1^{h_2-h_1-d_3}t_2^{h_3-h_2-d_4}
 \sum_{n_1+n_2+n_3\leq N}
 C^{(c)}_{n_1n_2n_3}
 \left(\frac{z}{t_1}\right)^{n_1}
 \left(\frac{t_1}{t_2}\right)^{n_2}t_2^{n_3},
 \label{eq:sixpoint-value-c-recursion}
\end{equation}
where $C^{(c)}$ is generated by the fixed-weight central-charge recursion.
The displayed curves use $N=10$ for the elliptic recursion and $N=20$ for
the $c$-recursion.  Their observed truncation scales are measured by
$8\to10$ and $18\to20$, respectively.

\begin{figure}[H]
 \centering
 \includegraphics[width=0.9\textwidth]{\detokenize{../../Data Set/h-Recursion/sphere_six_pillow_h_vs_c_fixed_t1_t2.png}}
 \caption{The sphere six-point chiral block as the first mobile insertion
 $z$ varies while $t_1,t_2$ and all weights remain fixed.  The upper-panel
 curves are visually indistinguishable.  The lower panel compares their
 relative difference with the independently observed elliptic and
 central-charge truncation shifts.}
 \label{fig:sixpoint-pillow-h-vs-c-fixed-t1-t2}
\end{figure}

Across all $61$ points, the median relative difference is
$9.25\times10^{-6}$ and the maximum is $2.91\times10^{-5}$, attained at
$z=0.20$.  At that endpoint the elliptic $8\to10$ shift is
$1.78\times10^{-6}$ and the $c$-recursion $18\to20$ shift is
$4.98\times10^{-5}$.  Point by point, the disagreement is smaller than the
sum of these two shifts; its maximum ratio to that sum is $0.563$.  The
six-point elliptic and central-charge recursions therefore agree throughout
the scan within an uncertainty estimated independently from each series.

\paragraph{Posterior convergence of the order-ten elliptic truncation.}
The $8\to10$ shift used in Fig.~\ref{fig:sixpoint-pillow-h-vs-c-fixed-t1-t2}
is deliberately conservative: it includes two complete shells already
retained in the order-ten answer.  To estimate the omitted tail more
directly, we extended only the elliptic recursion to total degree twelve.
For
\begin{equation}
 \epsilon_N(z):=
 \frac{|\cF_6^{(h,N)}(z)-\cF_6^{(h,N-1)}(z)|}
 {|\cF_6^{(h,N)}(z)|},
 \label{eq:sixpoint-shell-size}
\end{equation}
the shell sizes over the same $61$ points are
\begin{center}
\begin{tabular}{c|c|c}
shell & median $\epsilon_N$ & maximum $\epsilon_N$\\
\hline
$N=10$ & $4.91\times10^{-7}$ & $1.54\times10^{-6}$\\
$N=11$ & $2.26\times10^{-7}$ & $6.47\times10^{-7}$\\
$N=12$ & $1.07\times10^{-7}$ & $2.62\times10^{-7}$
\end{tabular}
\end{center}
There is mild pointwise alternation, so the most useful diagnostic is the
combined two-shell change
\begin{equation}
 \eta_{10}^{(12)}(z):=
 \frac{|\cF_6^{(h,12)}(z)-\cF_6^{(h,10)}(z)|}
 {|\cF_6^{(h,12)}(z)|}.
 \label{eq:sixpoint-order-ten-posterior}
\end{equation}
Its minimum, median, and maximum are respectively
$1.18\times10^{-7}$, $3.19\times10^{-7}$, and
$9.09\times10^{-7}$.  The maximum occurs at $z=0.015$; at the midpoint
$z=0.1075$ and the opposite endpoint $z=0.20$ the values are
$3.19\times10^{-7}$ and $1.18\times10^{-7}$.  Thus a conservative
posterior statement for this real-moduli scan is
\begin{equation}
 \boxed{\text{relative truncation uncertainty of the order-ten elliptic
 block is below }10^{-6}.}
 \label{eq:sixpoint-order-ten-convergence}
\end{equation}
This is an empirical next-shell estimate, not a rigorous tail bound.  It is
nevertheless much smaller than the maximum $2.90\times10^{-5}$ elliptic--$c$
difference and the maximum $4.98\times10^{-5}$ $c$-recursion
$18\to20$ shift.  Hence the value comparison is limited by the slower
central-charge series, not by the order-ten elliptic truncation.

The script, CSV table, JSON summary, and vector/raster plots are stored in
\path{Code/sphere_five_kummer_h_recursion/} and
\path{Data Set/h-Recursion/}.

The formulas above are chiral.  The full nonchiral block is obtained by
multiplying by the antiholomorphic counterpart with barred weights and
coordinates.

\end{document}
